%!TEX root = ../../thesis.tex
% generiert von thesis/scripts/make_appendix_prompts.py — nicht von Hand aendern
\subsection*{Phase 1: Business Understanding (\texttt{prompts.bu\_initial\_prompt})}
\begin{AppCPrompt}
You are an ML engineer working in the current directory. Write Python code to
\AppCCont@^files and run it with Bash. Use pandas and scikit-learn. Work exclusively in the
\AppCCont@^current directory; do not access anything outside it (e.g. via \AppCQ@^..\AppCQ@^).
There is no client or domain expert to consult: make and state your own
\AppCCont@^assumptions.
TASK (Business Understanding): Your business task is:
  "\AppCVal@Determine whether a person\AppCQ@^s annual income exceeds 50K dollars from census^
\AppCCont@^\AppCVal@attributes.^"
Inspect \AppCQ@^data/train.parquet\AppCQ@^ with pandas (columns, dtypes, example values) and
\AppCCont@^derive an ML task specification from it. Write EXACTLY ONE of the two variants
\AppCCont@^as JSON to \AppCQ@^bu/task_spec.json\AppCQ@^:
1) If the task can be solved with this data:
   {"task_type": "classification", "target_column": "<column>", "metric":
\AppCCont@^"f1"|"accuracy", "feature_columns": [<all remaining columns>],
\AppCCont@^"critical_columns": [<optional: columns that should be checked especially
\AppCCont@^carefully in the later data understanding (DU) phase, e.g. suspected ID or
\AppCCont@^leakage columns>]}
2) If no target column can be derived from the task and the data (the target
\AppCCont@^variable does not exist in the data or is ambiguous):
   {"cannot_determine": true, "reason": "<justification>", "missing_info":
\AppCCont@^[<what is missing>]}
Do not guess a target column that does not fit the task \AppCDash@^ an honest
\AppCCont@^cannot_determine is better than an invented target. Train NOTHING, only analyze.
Where the reference model asks for reports, documentation or plots, keep them
\AppCCont@^brief (e.g. a short notes file); they do not replace the required output files.
Follow the generic tasks of this phase from the CRISP-DM 1.0 reference model
\AppCCont@^(quoted verbatim):
\AppCCrisp@bu^
\end{AppCPrompt}
\subsection*{Phase 2: Data Understanding (\texttt{prompts.du\_initial\_prompt})}
\begin{AppCPrompt}
You are an ML engineer working in the current directory. Write Python code to
\AppCCont@^files and run it with Bash. Use pandas and scikit-learn. Work exclusively in the
\AppCCont@^current directory; do not access anything outside it (e.g. via \AppCQ@^..\AppCQ@^).
There is no client or domain expert to consult: make and state your own
\AppCCont@^assumptions.
Business task: "\AppCVal@Determine whether a person\AppCQ@^s annual income exceeds 50K dollars^
\AppCCont@^\AppCVal@from census attributes.^"
TASK (Data Understanding): Analyze \AppCQ@^data/train.parquet\AppCQ@^ READ-ONLY \AppCDash@^ do not
\AppCCont@^modify any data and do not create any new data files. Write a diagnostic report
\AppCCont@^as JSON to \AppCQ@^du/data_report.json\AppCQ@^ with EXACTLY these fields:
- "n_rows": <int> row count of train.parquet
- "columns": {"<column>": {"dtype": "numeric"|"bool"|"categorical", "n_missing":
\AppCCont@^<int>}} for EVERY column. dtype rule: boolean pandas dtype -> "bool", otherwise
\AppCCont@^numeric dtype -> "numeric", everything else -> "categorical".
- "class_balance": {"<str(class value)>": <int>} for the target column \AppCQ@^\AppCVal@col_14^\AppCQ@^
- "quality_issues": [<string>] \AppCDash@^ data quality findings in your own words (e.g.
\AppCCont@^missing values, suspicious columns); empty list if nothing stands out
- "excluded_columns": [{"column": "<column>", "reason": "<why>"}] \AppCDash@^ columns
\AppCCont@^that are NOT suitable as features (e.g. IDs, constants, suspected leakage
\AppCCont@^columns that separate the target suspiciously perfectly); empty list if none
Columns flagged in the business understanding (BU) phase \AppCDash@^ check them
\AppCCont@^especially carefully for reasons to exclude them: [\AppCQ@^\AppCVal@fnlwgt^\AppCQ@^, \AppCQ@^\AppCVal@education^\AppCQ@^,
\AppCCont@^\AppCQ@^\AppCVal@education-num^\AppCQ@^].
If the data fundamentally contradicts the task (e.g. the column \AppCQ@^\AppCVal@col_14^\AppCQ@^ does
\AppCCont@^not match the business task or is not a sensible target for it), instead write
\AppCCont@^ONLY {"spec_mismatch": true, "spec_mismatch_reason": "<why>"}.
IMPORTANT: Compute all numbers exactly with pandas \AppCDash@^ no estimates. Only
\AppCCont@^analyze, do not train anything.
Where the reference model asks for reports, documentation or plots, keep them
\AppCCont@^brief (e.g. a short notes file); they do not replace the required output files.
Follow the generic tasks of this phase from the CRISP-DM 1.0 reference model
\AppCCont@^(quoted verbatim):
\AppCCrisp@du^
\end{AppCPrompt}
\subsection*{Phase 3: Data Preparation (\texttt{prompts.initial\_prompt})}
\begin{AppCPrompt}
You are an ML engineer working in the current directory. Write Python code to
\AppCCont@^files and run it with Bash. Use pandas and scikit-learn. Work exclusively in the
\AppCCont@^current directory; do not access anything outside it (e.g. via \AppCQ@^..\AppCQ@^).
There is no client or domain expert to consult: make and state your own
\AppCCont@^assumptions.
Business task: "\AppCVal@Determine whether a person\AppCQ@^s annual income exceeds 50K dollars^
\AppCCont@^\AppCVal@from census attributes.^"
TASK (Data Preparation): Read \AppCQ@^data/train.parquet\AppCQ@^ (feature columns + column
\AppCCont@^\AppCQ@^\AppCVal@col_14^\AppCQ@^). Prepare the data (handle missing values, make all features numeric \AppCDash@^
\AppCCont@^int/float; convert booleans to 0/1). Write the result to \AppCQ@^prep/prepared.parquet\AppCQ@^
\AppCCont@^\AppCDash@^ feature columns + column \AppCQ@^\AppCVal@col_14^\AppCQ@^ with its original values, no NaNs. Save
\AppCCont@^your preprocessing code in \AppCQ@^prep/\AppCQ@^ so that it can be applied identically to new
\AppCCont@^data. Use only \AppCQ@^data/train.parquet\AppCQ@^; the test labels are not available.
Where the reference model asks for reports, documentation or plots, keep them
\AppCCont@^brief (e.g. a short notes file); they do not replace the required output files.
Follow the generic tasks of this phase from the CRISP-DM 1.0 reference model
\AppCCont@^(quoted verbatim):
\AppCCrisp@prep^
FINDINGS from the data understanding (DU) phase (take them into account):
- \AppCVal@Missing values present in three categorical columns, stored as NaN (originally^
\AppCCont@^\AppCVal@the Adult dataset\AppCQ@^s \AppCQ@^?\AppCQ@^ token): workclass=1970 (5.8%), occupation=1979 (5.8%),^
\AppCCont@^\AppCVal@native-country=587 (1.7%). workclass and occupation are missing on an almost^
\AppCCont@^\AppCVal@identical set of rows.^
- \AppCVal@education and education-num encode the same information: a perfect 1:1^
\AppCCont@^\AppCVal@bijection (each education level maps to exactly one education-num 1..16 and vice^
\AppCCont@^\AppCVal@versa). Keeping both is redundant and introduces collinearity.^
- \AppCVal@fnlwgt is a census sampling/population weight, not a personal attribute of the^
\AppCCont@^\AppCVal@individual; it is ~66% unique (near-continuous) and not a sensible income^
\AppCCont@^\AppCVal@predictor.^
- \AppCVal@capital-gain (91.9% zeros) and capital-loss (95.2% zeros) are highly^
\AppCCont@^\AppCVal@zero-inflated and right-skewed; capital-gain has a sentinel-like max of 99999.^
- \AppCVal@Target col_14 is imbalanced: ~23.9% >50K vs ~76.1% <=50K, so accuracy would be^
\AppCCont@^\AppCVal@misleading (prefer F1 on the positive class).^
Do NOT use these columns as features (excluded in the data understanding (DU)
\AppCCont@^phase):
- \AppCVal@fnlwgt^: \AppCVal@Census final sampling weight indicating how many people the record^
\AppCCont@^\AppCVal@represents in the population, not a personal attribute; near-continuous (66%^
\AppCCont@^\AppCVal@unique) and not predictive of an individual\AppCQ@^s income. Unsuitable as a feature.^
- \AppCVal@education^: \AppCVal@Redundant with education-num: perfect 1:1 mapping between the two.^
\AppCCont@^\AppCVal@Dropping the categorical duplicate and keeping the ordinal numeric education-num^
\AppCCont@^\AppCVal@avoids collinearity/duplication.^
\end{AppCPrompt}
\subsection*{Phase 4: Modeling (\texttt{prompts.initial\_prompt})}
\begin{AppCPrompt}
You are an ML engineer working in the current directory. Write Python code to
\AppCCont@^files and run it with Bash. Use pandas and scikit-learn. Work exclusively in the
\AppCCont@^current directory; do not access anything outside it (e.g. via \AppCQ@^..\AppCQ@^).
There is no client or domain expert to consult: make and state your own
\AppCCont@^assumptions.
Business task: "\AppCVal@Determine whether a person\AppCQ@^s annual income exceeds 50K dollars^
\AppCCont@^\AppCVal@from census attributes.^"
TASK (Modeling): Read \AppCQ@^prep/prepared.parquet\AppCQ@^ (feature columns + column
\AppCCont@^\AppCQ@^\AppCVal@col_14^\AppCQ@^). Save the best trained model according to the metric in
\AppCCont@^\AppCQ@^bu/task_spec.json\AppCQ@^ with pickle to \AppCQ@^model/model.pkl\AppCQ@^ and write
\AppCCont@^\AppCQ@^model/metrics.json\AppCQ@^ with {"metric": "<the metric from \AppCQ@^bu/task_spec.json\AppCQ@^>",
\AppCCont@^"val_score": <float>, "dummy_score": <float>} \AppCDash@^ val_score: this metric for the
\AppCCont@^saved model on the validation data of your test design (a held-out validation
\AppCCont@^split or the cross-validation mean); dummy_score: the same metric for
\AppCCont@^DummyClassifier(strategy=\AppCQ@^most_frequent\AppCQ@^), measured the same way. Write
\AppCCont@^\AppCQ@^model/model.pkl\AppCQ@^ and \AppCQ@^model/metrics.json\AppCQ@^ as soon as you have a first working
\AppCCont@^model and overwrite them whenever you find a better one. The test labels are not
\AppCCont@^available; build your test design from \AppCQ@^prep/prepared.parquet\AppCQ@^ only.
Where the reference model asks for reports, documentation or plots, keep them
\AppCCont@^brief (e.g. a short notes file); they do not replace the required output files.
Follow the generic tasks of this phase from the CRISP-DM 1.0 reference model
\AppCCont@^(quoted verbatim):
\AppCCrisp@model^
\end{AppCPrompt}
\subsection*{Phase 5: Evaluation (\texttt{prompts.initial\_prompt})}
\begin{AppCPrompt}
You are an ML engineer working in the current directory. Write Python code to
\AppCCont@^files and run it with Bash. Use pandas and scikit-learn. Work exclusively in the
\AppCCont@^current directory; do not access anything outside it (e.g. via \AppCQ@^..\AppCQ@^).
There is no client or domain expert to consult: make and state your own
\AppCCont@^assumptions.
Business task: "\AppCVal@Determine whether a person\AppCQ@^s annual income exceeds 50K dollars^
\AppCCont@^\AppCVal@from census attributes.^"
TASK (Evaluation): Load the model from \AppCQ@^model/model.pkl\AppCQ@^ and the test features
\AppCCont@^from \AppCQ@^data/test_features.parquet\AppCQ@^ (contains NO target column). Apply exactly the
\AppCCont@^same preprocessing as in training (reuse the code in \AppCQ@^prep/\AppCQ@^) and produce a
\AppCCont@^prediction for every row. Write the predictions to \AppCQ@^eval/predictions.parquet\AppCQ@^ \AppCDash@^
\AppCCont@^exactly one column \AppCQ@^prediction\AppCQ@^ containing class labels in the same form as the
\AppCCont@^values of the target column in \AppCQ@^data/train.parquet\AppCQ@^ (not encoded codes or
\AppCCont@^probabilities), same row count and row order as test_features. Do not retrain or
\AppCCont@^replace \AppCQ@^model/model.pkl\AppCQ@^ in this phase; note any recommended further iterations
\AppCCont@^in your notes.
Where the reference model asks for reports, documentation or plots, keep them
\AppCCont@^brief (e.g. a short notes file); they do not replace the required output files.
Follow the generic tasks of this phase from the CRISP-DM 1.0 reference model
\AppCCont@^(quoted verbatim):
\AppCCrisp@eval^
\end{AppCPrompt}
